%! Equivalent Representations of Polynomial and Rational Expressions — Unit 1, Lesson 1.11 — Guided Notes
\documentclass[10pt]{article}
\usepackage{precalculus-article}
\usepackage{precalculus-boxes}
\newcommand{\TallMath}[1]{$\displaystyle #1\rule[-1.4em]{0pt}{3.2em}$}

\begin{document}

\pageheader{Unit 1, Lesson 1.11}{Guided Notes}
\namedateperiod

\vspace{0.08in}

\begin{objectivebox}
By the end of this lesson, I will be able to\ldots
\begin{enumerate}[itemsep=2pt]
  \item Rewrite a polynomial or rational expression in \blank{3cm} form or \blank{3cm} form
        and explain what each form reveals. \hfill\textit{(LO 1.11.A)}
  \item Perform polynomial \blank{3.5cm} to write $f(x) = g(x)q(x) + r(x)$ and use the result
        to find \blank{3cm} asymptotes. \hfill\textit{(LO 1.11.B)}
  \item Apply the \blank{3cm} theorem and Pascal's Triangle to expand $(a+b)^n$. \hfill\textit{(LO 1.11.C)}
\end{enumerate}
\end{objectivebox}

\vspace{0.06in}

\begin{vocabbox}
\termblanklong{factored form}
\termblanklong{standard form}
\termblanklong{polynomial long division}
\termblanklong{quotient}
\termblanklong{remainder}
\termblanklong{slant asymptote}
\termblanklong{binomial theorem}
\termblanklong{Pascal's Triangle}
\end{vocabbox}

\vspace{0.06in}

\begin{hookbox}
Three rational functions are written on the board:
\[
  f(x) = \frac{2x^2 - 8}{x^2 + x - 6}
  \qquad
  g(x) = \frac{2(x-2)(x+2)}{(x-2)(x+3)}
  \qquad
  h(x) = \frac{2(x+2)}{x+3}, \quad x \ne 2
\]
These are the \textit{same function} in three different faces. \blank{4cm} form reveals zeros and
holes. \blank{4cm} form reveals end behavior. The simplified form reveals the \blank{3cm} of the
restricted domain.
\end{hookbox}

\vspace{0.06in}

\begin{notesbox}{LO 1.11.A --- Factored Form vs.\ Standard Form}
\renewcommand{\arraystretch}{1.5}
\begin{tabularx}{\linewidth}{@{} >{\bfseries}p{0.22\linewidth} X X @{}}
 & \textbf{Factored Form} & \textbf{Standard Form} \\
\midrule
What it reveals & \blank{4.5cm} & \blank{4.5cm} \\[4pt]
                & zeros, holes, asymptotes & end behavior, degree, leading coefficient \\[4pt]
Example (poly)  & $(x-1)(x+3)(x-2)$ & $x^3 - 7x + 6$ (expand) \\[4pt]
Example (rational) & $\dfrac{(x-2)(x+2)}{(x-2)(x+3)}$ & $\dfrac{x^2-4}{x^2+x-6}$ \\
\end{tabularx}

\medskip
\textbf{Key idea (EK 1.11.A.1):} Factored form reveals \blank{3cm}, \blank{3cm},
\blank{3cm}, and domain restrictions.

\textbf{Key idea (EK 1.11.A.2):} Standard form reveals \blank{3.5cm}.
\end{notesbox}

\vspace{0.06in}

\begin{notesbox}{LO 1.11.B --- Polynomial Long Division}
\textbf{Division Algorithm:} If $f$ is divided by $g$ (with $\deg g \geq 1$), then
\[
  f(x) = g(x)\cdot q(x) + r(x), \quad \text{where } \deg(r) < \deg(g).
\]

\textbf{Steps:}
\begin{enumerate}[itemsep=3pt]
  \item Write numerator and denominator in \blank{3cm} order (fill gaps with $0$ coefficients).
  \item Divide the \blank{3cm} term of the dividend by the leading term of the divisor. Write this as the first term of $q(x)$.
  \item Multiply that term by the \blank{3cm} and subtract from the dividend.
  \item Bring down the next term and repeat until $\deg(\text{remainder}) < \deg(\text{divisor})$.
\end{enumerate}

\medskip
\textbf{Example:} Divide $f(x) = 2x^2 + 3x - 5$ by $g(x) = x - 1$.

\smallskip
\begin{tabular}{r|l}
\multicolumn{1}{r}{} & \blank{5cm} \\[2pt]
\cline{2-2}
$x - 1$ & $2x^2 + 3x - 5$ \\
         & $- (\blank{4cm})$ \\
\cline{2-2}
         & \blank{4cm} \\
         & $- (\blank{3cm})$ \\
\cline{2-2}
         & \blank{3cm} \\
\end{tabular}

\smallskip
Result: $q(x) = $ \blank{3cm}, $r(x) = $ \blank{2cm}, so $f(x) = (x-1)(\blank{3cm}) + \blank{2cm}$.

\medskip
\textbf{Slant Asymptote (EK 1.11.B.2):} When $\deg(f) = \deg(g) + 1$, the slant asymptote of
$\dfrac{f(x)}{g(x)}$ is $y = q(x)$ (the \blank{4cm} after long division).
\end{notesbox}

\vspace{0.06in}

\begin{notesbox}{LO 1.11.C --- Binomial Theorem and Pascal's Triangle}
\textbf{Pascal's Triangle} (fill in the missing entries):

\begin{center}
\begin{tabular}{ccccccccccc}
  &   &   &   & 1 &   &   &   &   \\ % n=0
  &   &   & 1 &   & 1 &   &   &   \\ % n=1
  &   & 1 &   & 2 &   & 1 &   &   \\ % n=2
  & 1 &   & 3 &   & 3 &   & 1 &   \\ % n=3
1 &   & \blank{0.7cm} &   & \blank{0.7cm} &   & \blank{0.7cm} &   & 1 \\ % n=4
\end{tabular}
\end{center}

\textbf{Rule:} Each entry equals the \blank{5cm} of the two entries above it.

\medskip
\textbf{Binomial Theorem (EK 1.11.C.1):}
\[
  (a + b)^n = \sum_{k=0}^{n} \binom{n}{k} a^{n-k} b^k
\]
The coefficients $\binom{n}{k}$ match row $n$ of Pascal's Triangle.

\medskip
\textbf{Example:} Expand $(x + 2)^3$.

Using row 3 of Pascal's Triangle ($1,\ 3,\ 3,\ 1$):
\[
  (x+2)^3 = 1\cdot x^3 + 3\cdot x^2(\blank{1cm}) + 3\cdot x(\blank{1.5cm}) + 1\cdot(\blank{1.5cm})
           = \blank{6cm}
\]
\end{notesbox}

\end{document}
